\documentclass[11pt]{article}

\usepackage[margin=1in]{geometry}
\usepackage{enumitem}
\usepackage{amsmath, amssymb}

\setlist[enumerate]{itemsep=0.6em, topsep=0.6em}

\begin{document}

\begin{center}
  {\Large \textbf{CS 225 Examination}}\\[0.25em]
  {\large \textbf{Data Structures and Algorithms in C++ -- Answer Key}}\\[0.5em]
\end{center}

\begin{enumerate}[label=\textbf{\arabic*.}]

  \item \textbf{(A)} $O(\log n)$.\\
  In a balanced BST (like AVL or Red-Black tree), the height is $O(\log n)$. Each comparison 
  eliminates roughly half the remaining nodes, leading to logarithmic search time.

  \item \textbf{(A)} A shallow copy copies pointer values; a deep copy allocates new memory and 
  copies the data.\\
  Shallow copy duplicates pointers, so both objects point to the same memory. Deep copy creates 
  new memory and copies actual data, ensuring independence. This is critical for classes with 
  dynamically allocated members.

  \item \textbf{(C)} Stack.\\
  Function calls follow LIFO (Last In, First Out) semantics: the most recently called function 
  returns first. A stack naturally models this with push for function entry and pop for return.

  \item \textbf{(B)} $O(n^2)$.\\
  When the pivot is consistently the smallest or largest element (e.g., already sorted array with 
  first element as pivot), each partition reduces size by only 1, leading to $n + (n-1) + \cdots + 1 
  = O(n^2)$ comparisons.

  \item \emph{Sample answer:}\\
  \textbf{Amortized analysis}: Analyzes the average cost per operation over a sequence of operations, 
  rather than worst-case for a single operation.

  \textbf{Dynamic array example}:
  \begin{itemize}
    \item When full, the array doubles in size, requiring copying $n$ elements: $O(n)$ for that insertion
    \item However, this happens only at sizes $1, 2, 4, 8, 16, \ldots$, i.e., every $n$ insertions
    \item Total cost for $n$ insertions: $n + (1 + 2 + 4 + \cdots + n) < n + 2n = 3n$
    \item Average (amortized) cost per insertion: $3n/n = O(1)$
  \end{itemize}

  The occasional expensive resize is ``amortized'' over many cheap insertions, yielding constant 
  average time.

  \item \emph{Sample answer:}\\
  \textbf{Min-heap}: A binary tree where each parent is smaller than its children. Provides $O(\log n)$ 
  insertion and $O(1)$ access to minimum element, $O(\log n)$ deletion of minimum.

  \textbf{Priority queue}: An abstract data type (ADT) that supports insertion and extraction of 
  the highest-priority element. A min-heap is a common implementation.

  \textbf{Dijkstra's algorithm with priority queue}:
  \begin{itemize}
    \item Use min-heap to always extract the vertex with smallest tentative distance
    \item Insert/update: $O(\log V)$ per operation
    \item Each vertex inserted/extracted once, each edge causes potential decrease-key
    \item Time complexity: $O((V + E) \log V)$ with binary heap, or $O(V \log V + E)$ with Fibonacci heap
  \end{itemize}

  \item \emph{Sample answer:}\\
  \textbf{Rule of Three}: If a class manages resources (e.g., dynamic memory), define:
  \begin{enumerate}
    \item Destructor (to free resources)
    \item Copy constructor (for deep copying)
    \item Copy assignment operator (for deep assignment)
  \end{enumerate}
  This prevents double-deletion and dangling pointers.

  \textbf{Rule of Five} (C++11): Also define:
  \begin{enumerate}
    \setcounter{enumi}{3}
    \item Move constructor (for efficient transfer of resources)
    \item Move assignment operator
  \end{enumerate}
  This enables move semantics for performance optimization.

  \textbf{Rule of Zero}: Prefer using RAII classes (std::unique\_ptr, std::vector) that manage 
  resources automatically. If no explicit resource management is needed, don't define any of the 
  special member functions—let the compiler generate them. This is the safest and most maintainable 
  approach.

  \item \emph{Sample answer:}\\
  \textbf{std::vector}:
  \begin{itemize}
    \item \textbf{Advantages}: Contiguous memory (cache-friendly), fast random access $O(1)$, 
    good iteration performance
    \item \textbf{Disadvantages}: Insertion/deletion in middle is $O(n)$, reallocation on growth
    \item \textbf{Use when}: Random access needed, sequential iteration, append-heavy workload
  \end{itemize}

  \textbf{std::list} (doubly-linked list):
  \begin{itemize}
    \item \textbf{Advantages}: $O(1)$ insertion/deletion at any position (if iterator available), 
    no reallocation, stable iterators
    \item \textbf{Disadvantages}: No random access ($O(n)$ to reach element), poor cache locality, 
    higher memory overhead (pointers per node)
    \item \textbf{Use when}: Frequent insertions/deletions in middle, need iterator stability, 
    splicing operations
  \end{itemize}

  \textbf{General guideline}: Default to std::vector due to cache efficiency. Use std::list only 
  when insertion/deletion patterns specifically benefit from $O(1)$ operations at known positions.

\end{enumerate}

\end{document}
